% !TEX root = ../main.tex

\chapter*{致\quad 谢}
\addcontentsline{toc}{chapter}{致\quad 谢}

行文至此，
我的本科毕业论文也接近尾声。

回顾整个毕业设计过程，
从最初确定选题，
到完成系统需求分析、架构设计、功能开发、识别实验和论文撰写，
这段经历既漫长又充实。

在 HearBridge 系统的设计与实现过程中，
我不仅进一步熟悉了前后端分离、移动端开发、后台管理系统、机器学习服务和多端协同架构等工程实践，
也更加深刻地认识到，
一个完整系统的实现并不只是完成某一个页面或某一个算法，
而是需要在需求、数据、模型、接口、部署、测试和用户体验之间不断权衡。

首先，
我要感谢我的指导老师。

在毕业设计选题、系统设计、论文结构和写作规范等方面，
老师给予了我重要的指导和帮助，
使我能够逐步明确研究方向，
并将项目内容整理为较为完整的毕业论文。

在论文撰写过程中，
老师对论文结构、章节安排和表述方式提出了宝贵意见，
帮助我不断修改和完善本文。

其次，
我要感谢武汉大学计算机学院在本科阶段提供的学习环境和课程训练。

本科四年的学习让我逐渐建立起对计算机系统、软件工程、数据库、网络、算法和人工智能相关技术的基本认识，
也为本次毕业设计中多端系统开发和手势识别服务实现奠定了基础。

在完成 HearBridge 系统的过程中，
我切身体会到课堂知识与工程实践之间的联系，
也认识到自己仍有许多需要继续学习和提升的地方。

同时，
我要感谢在项目开发和论文撰写过程中给予我帮助的同学和朋友。

在系统调试、实验分析、论文格式整理和问题排查过程中，
他们给了我很多支持与鼓励。

这些交流帮助我在遇到困难时及时调整思路，
也让我能够坚持完成整个毕业设计。

此外，
我要感谢开源社区提供的技术资源。

本文系统开发过程中使用了 Spring Boot、Vue、Element Plus、HarmonyOS ArkTS、MediaPipe、TensorFlow、FastAPI、MySQL、Redis、MinIO、JMeter、LaTeX 等工具和框架。

这些开源技术为 HearBridge 系统的实现提供了重要支撑，
也让我能够在有限时间内完成一个包含移动端、后端、管理端和识别服务的多端协同系统。

最后，
我要感谢一直支持我的家人。

他们在我的学习和生活中给予了我持续的理解与支持，
让我能够较为安心地完成本科阶段的学习与毕业设计。

毕业设计的完成并不意味着学习的结束。

通过本次项目，
我更加清楚地认识到自己在工程能力、模型训练、系统设计和论文表达方面仍然存在不足。

未来我会继续保持学习和实践，
在已有基础上不断提高自己的专业能力。

谨以此文，
向所有在本科阶段给予我帮助、支持和鼓励的人致以诚挚的感谢。